% Options for packages loaded elsewhere
\PassOptionsToPackage{unicode}{hyperref}
\PassOptionsToPackage{hyphens}{url}
\PassOptionsToPackage{dvipsnames,svgnames,x11names}{xcolor}
\documentclass[
  english,
  10pt,
  a4paper,
]{article}
\usepackage{xcolor}
\usepackage[top=22mm,bottom=20mm,left=18mm,right=18mm]{geometry}
\usepackage{amsmath,amssymb}
\setcounter{secnumdepth}{-\maxdimen} % remove section numbering
\usepackage{iftex}
\ifPDFTeX
  \usepackage[T1]{fontenc}
  \usepackage[utf8]{inputenc}
  \usepackage{textcomp} % provide euro and other symbols
\else % if luatex or xetex
  \usepackage{unicode-math} % this also loads fontspec
  \defaultfontfeatures{Scale=MatchLowercase}
  \defaultfontfeatures[\rmfamily]{Ligatures=TeX,Scale=1}
\fi
\usepackage{lmodern}
\ifPDFTeX\else
  % xetex/luatex font selection
\fi
% Use upquote if available, for straight quotes in verbatim environments
\IfFileExists{upquote.sty}{\usepackage{upquote}}{}
\IfFileExists{microtype.sty}{% use microtype if available
  \usepackage[]{microtype}
  \UseMicrotypeSet[protrusion]{basicmath} % disable protrusion for tt fonts
}{}
\makeatletter
\@ifundefined{KOMAClassName}{% if non-KOMA class
  \IfFileExists{parskip.sty}{%
    \usepackage{parskip}
  }{% else
    \setlength{\parindent}{0pt}
    \setlength{\parskip}{6pt plus 2pt minus 1pt}}
}{% if KOMA class
  \KOMAoptions{parskip=half}}
\makeatother
\usepackage{longtable,booktabs,array}
\usepackage{caption}
\captionsetup[table]{skip=6pt}
\newcounter{none} % for unnumbered tables
\usepackage{calc} % for calculating minipage widths
% Correct order of tables after \paragraph or \subparagraph
\usepackage{etoolbox}
\makeatletter
\patchcmd\longtable{\par}{\if@noskipsec\mbox{}\fi\par}{}{}
\makeatother
\ifLuaTeX
\usepackage[bidi=basic,shorthands=off]{babel}
\else
\usepackage[bidi=default,shorthands=off]{babel}
\fi
\ifLuaTeX
  \usepackage{selnolig} % disable illegal ligatures
\fi
\setlength{\emergencystretch}{3em} % prevent overfull lines
\providecommand{\tightlist}{%
  \setlength{\itemsep}{0pt}\setlength{\parskip}{0pt}}
% ---- Çılgın Kurye document theme: same look as the CSE 396 project proposal ----
% Compile with XeLaTeX (also works on Overleaf). Fonts: TeX Gyre Termes (Times), DejaVu Sans Mono.
\usepackage{fontspec}
\setmainfont{texgyretermes}[Extension=.otf, UprightFont=*-regular, BoldFont=*-bold, ItalicFont=*-italic, BoldItalicFont=*-bolditalic]
\setsansfont{texgyreheros}[Extension=.otf, UprightFont=*-regular, BoldFont=*-bold, ItalicFont=*-italic, BoldItalicFont=*-bolditalic]
\setmonofont{DejaVuSansMono}[Extension=.ttf, UprightFont=*, BoldFont=*-Bold, ItalicFont=*-Oblique, BoldItalicFont=*-BoldOblique, Scale=0.80]
\setmathfont{texgyretermes-math.otf}

\usepackage{xcolor}
\definecolor{ckblue}{HTML}{1F3A5F}
\definecolor{cklight}{HTML}{E8EEF5}
\definecolor{ckgrey}{HTML}{F2F2F2}
\definecolor{ckline}{HTML}{555555}
\definecolor{shadecolor}{HTML}{F4F5F7}

% Headings as in the proposal: serif, bold; subsections bold italic. Numbers are part of the heading text.
\setcounter{secnumdepth}{0}
\usepackage{titlesec}
\titleformat{\section}{\rmfamily\Large\bfseries}{}{0pt}{}
\titleformat{\subsection}{\rmfamily\large\bfseries\itshape}{}{0pt}{}
\titleformat{\subsubsection}{\rmfamily\normalsize\bfseries}{}{0pt}{}
\titlespacing*{\section}{0pt}{12pt}{4pt}
\titlespacing*{\subsection}{0pt}{8pt}{3pt}

% Header and footer (the first page has only the footer, like the proposal)
\usepackage{fancyhdr,lastpage}
\pagestyle{fancy}
\fancyhf{}
\fancyhead[L]{\footnotesize\itshape CSE 396 · Group 8 · Çılgın Kurye}
\fancyhead[R]{\footnotesize\itshape\ckmodule}
\fancyfoot[C]{\footnotesize Page \thepage\ of \pageref*{LastPage}}
\renewcommand{\headrulewidth}{0.4pt}
\setlength{\headheight}{14pt}
\fancypagestyle{ckfirst}{\fancyhf{}\renewcommand{\headrulewidth}{0pt}\fancyfoot[C]{\footnotesize Page \thepage\ of \pageref*{LastPage}}}

% Inline code: breaks only after . : ( , / _ and spaces, never inside a name
\usepackage[obeyspaces,spaces]{url}
\AtBeginDocument{%
  \renewcommand{\UrlBreaks}{\do\.\do\:\do\(\do\,\do\/\do\_\do\|}%
  \renewcommand{\UrlBigBreaks}{}%
  \renewcommand{\UrlNoBreaks}{}%
  \urlstyle{tt}}

% Tables: boxed grid as in the proposal, smaller text
\usepackage{etoolbox}
\AtBeginEnvironment{longtable}{\footnotesize\renewcommand{\arraystretch}{1.2}}
\setlength{\LTpre}{6pt}
\setlength{\LTpost}{8pt}
\setlength{\arrayrulewidth}{0.4pt}

% Code blocks: long lines wrap; diagrams stay in one piece
\usepackage{framed}
\usepackage{fvextra}
\RecustomVerbatimEnvironment{verbatim}{Verbatim}{fontsize=\footnotesize, breaklines}
\makeatletter
\@ifundefined{Shaded}{}{\renewenvironment{Shaded}{\begin{snugshade}}{\end{snugshade}}}
\@ifundefined{Highlighting}{}{\RecustomVerbatimEnvironment{Highlighting}{Verbatim}{commandchars=\\\{\}, fontsize=\footnotesize, breaklines, breakanywhere}}
\makeatother

% Figures drawn with TikZ
\usepackage{tikz}
\usetikzlibrary{arrows.meta,positioning,fit,calc,backgrounds,shapes.geometric}
\tikzset{
  ck/.style={font=\footnotesize, >={Stealth[length=5pt]}, line width=0.6pt},
  box/.style={draw=ckline, fill=white, rounded corners=2pt, align=center, inner sep=4pt, minimum height=7mm},
  mod/.style={box, fill=ckgrey},
  this/.style={box, fill=cklight, draw=ckblue, line width=0.9pt},
  group/.style={draw=ckline, dashed, rounded corners=3pt, inner sep=6pt},
  lbl/.style={font=\scriptsize, align=center, inner sep=1.5pt},
  msg/.style={->, draw=ckline},
}
\usepackage{caption}
\captionsetup[figure]{font=small, labelfont=bf, skip=4pt}
\newcommand{\ckfigure}[2]{\begin{figure}[htbp]\centering #1\caption{#2}\end{figure}}

\usepackage{enumitem}
\setlist{itemsep=1pt, topsep=3pt}
\setlength{\parskip}{4pt}
\setlength{\parindent}{0pt}
\raggedbottom
% ---- end of theme ----
\newcommand{\ckmodule}{M7 ESP32 Software and Communication}
\usepackage{bookmark}
\IfFileExists{xurl.sty}{\usepackage{xurl}}{} % add URL line breaks if available
\urlstyle{same}
% fallback for those not using the hyperref driver hyperxmp:
\makeatletter
\@ifundefined{xmpquote}{\newcommand{\xmpquote}[1]{#1}}{}
\makeatother
\hypersetup{
  pdftitle={M7 ESP32 Software and Communication},
  pdflang={en},
  colorlinks=true,
  linkcolor={black},
  filecolor={Maroon},
  citecolor={Blue},
  urlcolor={black},
  pdfcreator={LaTeX}}

\title{M7 ESP32 Software and Communication --- Module Documentation}
\author{}
\date{}


% ---- colours and drawings ----
% ---- colours ----
\definecolor{srv}{HTML}{2E6FB7}\colorlet{srvL}{srv!15}
\definecolor{cli}{HTML}{2E9E5B}\colorlet{cliL}{cli!15}
\definecolor{adm}{HTML}{E08A1E}\colorlet{admL}{adm!18}
\definecolor{brd}{HTML}{7B4FB3}\colorlet{brdL}{brd!15}
\definecolor{tcpc}{HTML}{1F5FA8}
\definecolor{udpc}{HTML}{D9531E}
\definecolor{obs}{HTML}{D7263D}
\definecolor{tA}{HTML}{00A6D6}
\definecolor{tB}{HTML}{3BB273}
\definecolor{tC}{HTML}{F2C14E}
\definecolor{tD}{HTML}{9B5DE5}

% ---- styles ----
\tikzset{
  nb/.style={rounded corners=3pt, line width=0.9pt, align=center, inner sep=4pt, minimum height=9mm, font=\footnotesize},
  sv/.style={nb, draw=srv, fill=srvL},
  cl/.style={nb, draw=cli, fill=cliL},
  ad/.style={nb, draw=adm, fill=admL},
  bd/.style={nb, draw=brd, fill=brdL},
  hwb/.style={nb, draw=black!50, fill=black!5},
  tcpa/.style={draw=tcpc, line width=1.5pt},
  udpa/.style={draw=udpc, line width=1.5pt, dash pattern=on 5pt off 2.5pt},
  inna/.style={draw=black!55, line width=1pt},
  al/.style={font=\scriptsize, fill=white, inner sep=1.2pt, rounded corners=1pt},
  grp/.style={draw=black!35, dashed, rounded corners=5pt, inner sep=7pt},
  glbl/.style={font=\scriptsize\bfseries, text=black!65},
  hl/.style={line width=2.6pt},
  st/.style={nb, rounded corners=8pt, minimum width=24mm},
  co/.style={draw=black!60, line width=0.5pt},
  cot/.style={font=\scriptsize, align=left, text=black!80},
}

% ---- system map; argument = module to highlight (M1 ... M8, or none) ----
\newcommand{\systemmap}[1]{%
\tikzset{M1/.style={},M2/.style={},M3/.style={},M4/.style={},M5/.style={},M6/.style={},M7/.style={},M8/.style={},none/.style={}}%
\tikzset{#1/.style={hl}}%
\begin{tikzpicture}[ck]
\node[cl, M5, minimum width=27mm] (m5) at (0,3) {\textbf{M5}\\Player phones};
\node[ad, M6, minimum width=27mm] (m6) at (0,0.4) {\textbf{M6}\\Admin panel};
\node[sv, M2, minimum width=27mm, minimum height=38mm] (m2) at (4.6,1.7) {\textbf{M2}\\Network\\and sync};
\node[sv, M1, minimum width=27mm, minimum height=38mm] (m1) at (8.4,1.7) {\textbf{M1}\\Game server\\and rules};
\node[sv, M3, minimum width=25mm] (m3) at (12,2.8) {\textbf{M3}\\Physics};
\node[sv, M4, minimum width=25mm] (m4) at (12,0.6) {\textbf{M4}\\Traffic and AI};
\node[grp, fit=(m2)(m1)(m3)(m4)] (pc) {};
\node[glbl, anchor=south west] at (pc.north west) {MAIN COMPUTER};
\node[bd, M7, minimum width=27mm] (m7) at (4.6,-2.7) {\textbf{M7}\\ESP32 software};
\node[bd, M8, minimum width=27mm] (m8) at (8.4,-2.7) {\textbf{M8}\\LED engine};
\node[hwb, minimum width=25mm] (led) at (12,-2.7) {344 LEDs};
\node[hwb, minimum width=27mm] (sm) at (4.6,-4.4) {Sensor and\\4 motors};
\node[grp, fit=(m7)(m8)(led)(sm)] (bb) {};
\node[glbl, anchor=north east] at (bb.south east) {LED CITY BOARD};
\draw[tcpa, <->] (m5.east) -- node[al]{WebSocket} (m5.east -| m2.west);
\draw[tcpa, <->] (m6.east) -- node[al]{WebSocket} (m6.east -| m2.west);
\draw[tcpa, <->] ([xshift=-6mm]m7.north) -- node[al]{TCP} ([xshift=-6mm]m2.south);
\draw[udpa, <->] ([xshift=6mm]m7.north) -- node[al]{UDP} ([xshift=6mm]m2.south);
\draw[inna, <->] (m2) -- (m1);
\draw[inna, <->] (m1.east |- m3) -- (m3);
\draw[inna, <->] (m1.east |- m4) -- (m4);
\draw[inna, ->] (m7) -- node[al]{state} (m8);
\draw[inna, ->] (m8) -- (led);
\draw[inna, <->] (m7) -- (sm);
\draw[tcpa] (-1.3,-6.0) -- ++(0.9,0) node[right, font=\scriptsize]{TCP: every message arrives, in order};
\draw[udpa] (4.0,-6.0) -- ++(0.9,0) node[right, font=\scriptsize]{UDP: fast, the newest message wins};
\draw[inna] (9.3,-6.0) -- ++(0.9,0) node[right, font=\scriptsize]{inside one program or chip};
\end{tikzpicture}}

% ---- LED board seen from above (6 x 6 roads, 1 unit = 10 cm) ----
\newcommand{\boardbase}{%
\fill[black!88, rounded corners=3pt] (-0.5,-0.5) rectangle (5.5,5.5);
\foreach \i in {0,...,4}{\foreach \j in {0,...,4}{%
  \fill[black!60] (\i+0.13,\j+0.13) rectangle (\i+0.87,\j+0.87);
  \fill[black!22] (\i+0.24,\j+0.24) rectangle (\i+0.52,\j+0.52);
  \fill[black!30] (\i+0.56,\j+0.30) rectangle (\i+0.78,\j+0.62);
  \fill[black!38] (\i+0.28,\j+0.58) rectangle (\i+0.66,\j+0.78);}}
\foreach \i in {0,...,5}{%
  \draw[black, line width=4.5pt] (\i,-0.05) -- (\i,5.05);
  \draw[black, line width=4.5pt] (-0.05,\i) -- (5.05,\i);}
\foreach \i in {0,...,5}{\foreach \k in {0,...,30}{%
  \fill[black!55] (\i,\k/6) circle (0.025);
  \fill[black!55] (\k/6,\i) circle (0.025);}}
\foreach \x/\y/\c in {0/4/tA, 4/4/tB, 0/0/tC, 4/0/tD}{%
  \fill[\c!30!black!60] (\x+0.13,\y+0.13) rectangle (\x+0.87,\y+0.87);
  \fill[\c] (\x+0.3,\y+0.3) rectangle (\x+0.7,\y+0.7);
  \fill[white] (\x+0.4,\y+0.5) circle (0.05); \fill[white] (\x+0.6,\y+0.5) circle (0.05);}}
\newcommand{\nmark}[3]{\node[circle, fill=yellow!85!orange, draw=black, line width=0.4pt, inner sep=0.6pt, minimum size=3.6mm, font=\tiny\bfseries] at (#1,#2) {#3};}
\newcommand{\bike}[4]{%
\fill[#3!35] ($(#1,#2)+(#4:0.40)$) circle (0.05);
\fill[#3!65] ($(#1,#2)+(#4:0.21)$) circle (0.075);
\fill[#3, draw=white, line width=0.4pt] (#1,#2) circle (0.12);}
\newcommand{\dpoint}[3]{\draw[#3!50, line width=0.8pt] (#1,#2) circle (0.26); \draw[#3, line width=1.4pt] (#1,#2) circle (0.16); \fill[#3] (#1,#2) circle (0.07);}
\newcommand{\heavy}[2]{\draw[obs!60!black, line width=4pt] (#1) -- (#2); \draw[obs, line width=4pt, dash pattern=on 3pt off 5pt] (#1) -- (#2);}
\newcommand{\roadwork}[2]{\draw[obs, line width=4pt] (#1) -- (#2);}
\newcommand{\accident}[2]{\draw[obs, line width=4pt] (#1) -- (#2); \draw[white, line width=4pt, dash pattern=on 2.5pt off 3.5pt] (#1) -- (#2);}
\newcommand{\puddle}[2]{\draw[cyan!35, line width=4pt] (#1) -- (#2); \draw[white, line width=1.6pt, dash pattern=on 1pt off 2.5pt] (#1) -- (#2);}
\newcommand{\findme}[2]{\fill[white, opacity=0.25] (#1,#2) circle (0.36); \fill[white, opacity=0.5] (#1,#2) circle (0.24); \fill[white] (#1,#2) circle (0.13);}

% ---- phone game screen (7 x 3.4 units) ----
\newcommand{\phonemock}{%
\fill[black!85, rounded corners=9pt] (-0.3,-0.3) rectangle (7.3,3.7);
\fill[cyan!20] (0,1.7) rectangle (7,3.4);
\fill[cli!14] (0,0) rectangle (7,1.7);
\filldraw[fill=srv!30, draw=srv!70!black] (0.2,1.5) rectangle (1.0,3.0);
\filldraw[fill=adm!35, draw=adm!70!black] (1.1,1.6) rectangle (1.8,2.4);
\filldraw[fill=brd!30, draw=brd!70!black] (1.9,1.65) -- (2.6,1.65) -- (2.6,2.2) -- (2.25,2.5) -- (1.9,2.2) -- cycle;
\filldraw[fill=adm!35, draw=adm!70!black] (4.5,1.65) rectangle (5.1,2.3);
\filldraw[fill=srv!30, draw=srv!70!black] (5.2,1.55) rectangle (6.0,3.1);
\filldraw[fill=cli!35, draw=cli!70!black] (6.1,1.6) rectangle (6.85,2.6);
\fill[brown!60!black] (4.17,1.6) rectangle (4.23,1.8); \fill[cli!70!black] (4.2,1.95) circle (0.18);
\fill[black!55] (1.8,0) -- (5.2,0) -- (3.75,1.7) -- (3.25,1.7) -- cycle;
\draw[white, line width=0.8pt, dash pattern=on 4pt off 4pt] (3.5,0) -- (3.5,1.7);
\filldraw[fill=black!15, draw=black!70] (3.62,1.12) rectangle (3.88,1.32);
\fill[obs] (3.05,1.22) -- (3.15,1.48) -- (3.25,1.22) -- cycle;
\fill[black!85] (3.43,0.08) rectangle (3.57,0.3);
\fill[tA] (3.3,0.28) rectangle (3.7,0.6);
\fill[brown!65] (3.34,0.6) rectangle (3.66,0.82);
\fill[tA!60!black] (3.5,0.95) circle (0.13);
\node[fill=black!70, text=white, rounded corners=2pt, font=\scriptsize\bfseries, inner sep=2pt] at (3.5,3.15) {2:41};
\foreach \c/\s/\x in {tA/120/0.35,tB/95/0.95,tC/140/1.55,tD/80/2.15}
  \node[fill=\c, text=white, font=\tiny\bfseries, rounded corners=1pt, inner sep=1.5pt, minimum width=5mm] at (\x,3.15) {\s};
\fill[white, opacity=0.9] (6.45,2.85) circle (0.38); \draw[black!60] (6.45,2.85) circle (0.38);
\draw[-{Stealth[length=5pt]}, line width=1.8pt, obs] (6.45,2.85) ++(230:0.2) -- ++(50:0.5);
\fill[white, opacity=0.85] (0.55,0.5) circle (0.32); \fill[black!70] (0.43,0.5) -- (0.65,0.63) -- (0.65,0.37) -- cycle;
\fill[white, opacity=0.85] (1.4,0.5) circle (0.32); \fill[black!70] (1.52,0.5) -- (1.3,0.63) -- (1.3,0.37) -- cycle;
\node[fill=obs!80, text=white, rounded corners=3pt, font=\tiny\bfseries, minimum width=8mm, minimum height=6mm] at (5.6,0.5) {Brake};
\node[fill=cli, text=white, rounded corners=3pt, font=\tiny\bfseries, minimum width=8mm, minimum height=8mm] at (6.5,0.6) {Gas};
\node[fill=tC, text=black, rounded corners=3pt, font=\tiny\bfseries, minimum width=9mm, minimum height=5mm] at (6.15,1.45) {Deliver};
\node[fill=white, text=black, rounded corners=3pt, font=\tiny\bfseries, minimum width=9mm, minimum height=5mm] at (0.9,1.45) {Find me};}

% ---- side cut of the board (14 units wide) ----
\newcommand{\boardside}{%
\fill[brown!20] (-0.6,-0.35) rectangle (14.6,0);
\foreach \x/\h/\c in {1.0/0.7/black!18, 2.2/1.1/black!22, 4.6/1.9/black!25, 5.8/2.6/black!28, 7.4/2.2/black!25, 8.6/1.6/black!22, 10.9/1.0/black!18}
  \filldraw[fill=\c, draw=black!45] (\x,1.5) rectangle (\x+0.9,1.5+\h);
\filldraw[fill=tB!35, draw=tB!70!black] (12.4,1.5) rectangle (13.4,2.3);
\fill[black!70, rounded corners=1pt] (0.5,0) rectangle (1.1,0.9);
\fill[black!70, rounded corners=1pt] (12.9,0) rectangle (13.5,0.9);
\filldraw[fill=brown!45, draw=brown!70!black] (0,0.9) rectangle (14,1.3);
\fill[black!80] (0,1.3) rectangle (14,1.38);
\foreach \x in {0.3,0.9,...,13.8} \fill[black!35] (\x,1.42) circle (0.06);
\foreach \x/\c in {3.3/tA, 3.9/tA, 6.9/obs, 7.5/obs, 10.5/tC} \fill[\c] (\x,1.42) circle (0.07);
\fill[white, opacity=0.55] (0,1.38) rectangle (14,1.52);
\filldraw[fill=brd!40, draw=brd] (2.3,0.35) rectangle (4.3,0.85);
\node[font=\tiny] at (3.3,0.6) {ESP32};
\filldraw[fill=brd, draw=brd!60!black] (6.75,0.62) rectangle (7.25,0.88);
\filldraw[fill=adm!60, draw=adm!70!black] (1.8,0.6) circle (0.24);
\filldraw[fill=adm!60, draw=adm!70!black] (12.2,0.6) circle (0.24);
\filldraw[fill=black!15, draw=black!50] (9.3,0.35) rectangle (10.9,0.85);
\node[font=\tiny] at (10.1,0.6) {Power};}

\begin{document}

\thispagestyle{ckfirst}
\begin{center}
{\bfseries GEBZE TECHNICAL UNIVERSITY\par}
{Computer Engineering Department\par}
{CSE 396 Module Documentation\par}
\vspace{8pt}
{\LARGE\bfseries Çılgın Kurye\par}
\vspace{3pt}
{\large\itshape Module M7: ESP32 Software and Communication\par}
\vspace{6pt}
\textbf{Instructor:} Erkan Zergeroğlu\hspace{1.4em}\textbf{Company:} Board\&Byte\hspace{1.4em}\textbf{Group:} Group-8\hspace{1.4em}\textbf{Date:} 6 November 2026
\end{center}
\begin{center}\footnotesize\renewcommand{\arraystretch}{1.15}
\begin{tabular}{|p{0.27\linewidth}|p{0.15\linewidth}|p{0.27\linewidth}|p{0.15\linewidth}|}\hline
\textbf{Name} & \textbf{Student No.} & \textbf{Name} & \textbf{Student No.} \\ \hline
Tuğberk Akbulut & 230104004070 & Hasan Alperen Elmas & 1801042604 \\ \hline
Hasan Hüseyin Özdemir & 230104004048 & Muhib Taha Boy & 230104004165 \\ \hline
\end{tabular}
\end{center}
\vspace{2pt}





% ---- table and list helpers ----
\newcolumntype{L}[1]{>{\raggedright\arraybackslash}p{\dimexpr#1\linewidth-2\tabcolsep-\arrayrulewidth\relax}}
\newcolumntype{C}[1]{>{\centering\arraybackslash}p{\dimexpr#1\linewidth-2\tabcolsep-\arrayrulewidth\relax}}
\newcommand{\ckdone}{\item[$\boxtimes$]}
\newcommand{\ckopen}{\item[$\square$]}
\newcommand{\mdot}{$\bullet$}
\newcounter{msec}
\newcounter{chap}
% ---- end of helpers ----
\newcommand{\msection}[1]{\stepcounter{msec}\section{\arabic{msec}. #1}}

\msection{Purpose}

M7 is the software of the ESP32 inside the LED board, and the link between the board and the main
computer. It gets the game state over Wi-Fi and gives it to the LED engine (M8). It reads the motion
sensor and finds the strength and the place of a hit on the board. It drives the vibration motors.
It also reports the health of the board and the sensor signal to the admin.

\begin{longtable}[]{|L{0.18}|L{0.82}|}
\hline Runs on & ESP32 in the LED board: core 0 (the LED engine M8 runs on core 1) \\ \hline
Technology & C++ with the Arduino framework for ESP32 and FreeRTOS tasks; TCP and UDP over Wi-Fi with MessagePack; 6-axis motion sensor (accelerometer and gyroscope); 4 vibration motors with transistor drivers \\ \hline
Requirements & R7 R8 R10 R11 \\ \hline
\end{longtable}

\begin{figure}[htbp]\centering
\systemmap{M7}
\caption{Where M7 sits in the system. M7 talks to the server on two channels: TCP and UDP.}
\end{figure}

\msection{Requirements covered}

\begin{longtable}[]{|L{0.07}|L{0.25}|L{0.68}|}
\hline \textbf{ID} & \textbf{Requirement} & \textbf{What M7 does for it} \\ \hline \endhead
R7 & Earthquake & Finds hand hits and motor shaking with the sensor; strength 0--10 and a 3 × 3 zone; drives the motors for the admin button \\ \hline
R8 & Physical LED board & Brings the game state to the board over Wi-Fi; one ESP32 for the LEDs, the sensor and the motors \\ \hline
R10 & Network and measurement & Board delay measurement; sensor samples for the seismic graph; comes back after a lost link \\ \hline
R11 & Independent development & PC emulators for both sides of the link \\ \hline
\end{longtable}

\msection{Inputs and outputs}

\begin{longtable}[]{|L{0.08}|L{0.20}|L{0.18}|L{0.54}|}
\hline \textbf{Dir.} & \textbf{Data} & \textbf{From / to} & \textbf{Content} \\ \hline \endhead
In & Board state & M2, UDP & 10 to 30 per second: phase, players (team, tone, position, flags), delivery points, obstacles, earthquake \\ \hline
In & Config & M2, TCP & Player colours, brightness 0.1--1, LED current limit \\ \hline
In & Vibrate & M2, TCP & Motors to use, strength 0--1, duration of at most 5 s \\ \hline
In & Sensor & Hardware & 1000 samples per second \\ \hline
Out & Newest state & M8 & Latest board state, config and link state, on the same chip \\ \hline
Out & Earthquake event & M2, TCP, to M1 and M6 & Source (hand or motors), strength 0--10, centre, zone 0--8, confidence \\ \hline
Out & Sensor samples & M2, UDP, to M6 & 100 per second for the seismic graph, sent in groups of 10 \\ \hline
Out & Status & M2, TCP, to M6 & Once per second: frame rate, Wi-Fi signal, LED current, errors \\ \hline
\end{longtable}

\msection{Main components}

\begin{longtable}[]{|L{0.20}|L{0.52}|L{0.28}|}
\hline \textbf{Component} & \textbf{What it does} & \textbf{Related modules} \\ \hline \endhead
Network link & Joins Wi-Fi, finds the server, keeps the TCP channel and the UDP stream, connects again after a loss & M2 \\ \hline
Message codec & MessagePack encode and decode; checks the version and the order numbers; drops old or broken messages & M2 \\ \hline
State hand-over & Two-slot buffer: the LED engine always takes the newest full state and never waits & M8 \\ \hline
Hit detector & Finds hits in the sensor signal; computes the strength and the zone & M1, M6 \\ \hline
Motor controller & Runs vibrate commands with a soft start and safety limits & M1, M6 \\ \hline
Status and settings & Status reports and delay replies; Wi-Fi and calibration settings in flash, changed over USB & M6 \\ \hline
\end{longtable}

\msection{How it works}

\textbf{Two channels.} The board keeps one TCP channel and one UDP stream to the server. Messages that
must arrive go on TCP: the config, the vibrate commands, the earthquake events and the status
report. Data that comes again and again goes on UDP: the board state, the sensor samples and the
delay probes. If one UDP packet is lost, the next one replaces it, so the LEDs never wait. Every
message has an order number, so the board can drop a UDP packet that comes late and out of order.

\textbf{Two cores.} The ESP32 has two cores. Core 0 runs all M7 work: Wi-Fi, the network tasks, the
sensor, the hit detector and the motors. Core 1 runs only the LED engine of M8, 60 frames per
second. So Wi-Fi work never breaks the LED timing. The two sides share the state through a
two-slot buffer. M7 always writes into the free slot. When the slot is full, M7 switches one index
in one step. The LED engine reads the slot that the index shows, so it always gets a complete state
and it never waits for a lock.

\begin{figure}[htbp]\centering
\begin{tikzpicture}[ck]
\draw[brd, line width=1.2pt, rounded corners=6pt, fill=brd!4] (-0.4,-3.3) rectangle (13.2,1.5);
\node[font=\scriptsize\bfseries, text=brd, anchor=north west] at (-0.3,1.45) {ESP32};
\node[font=\scriptsize\bfseries, text=brd] at (3.3,1.1) {Core 0: M7};
\node[font=\scriptsize\bfseries, text=brd] at (10.6,1.1) {Core 1: M8};
\node[bd, minimum width=22mm] (wifi) at (1.2,0.2) {Wi-Fi and\\network};
\node[bd, minimum width=22mm] (cod) at (4.4,0.2) {Message\\codec};
\node[bd, minimum width=22mm] (sen) at (1.2,-2.0) {Sensor\\1000 per s};
\node[bd, minimum width=22mm] (hit) at (4.4,-2.0) {Hit\\detector};
\node[bd, minimum width=16mm] (mot) at (7.0,-2.0) {Motor\\control};
\node[nb, draw=adm, fill=admL, minimum width=18mm, minimum height=14mm] (buf) at (7.4,0.2) {Two-slot\\buffer};
\node[nb, draw=brd, fill=brd!25, minimum width=22mm, minimum height=24mm, line width=1.4pt] (led) at (10.6,-0.7) {LED engine\\60 frames\\per second};
\draw[inna,->] (wifi) -- (cod);
\draw[inna,->] (cod) -- node[al]{state} (buf);
\draw[inna,->] (buf) -- (led.west |- buf);
\draw[inna,->] (sen) -- (hit);
\draw[inna,->] (hit.north) -- node[al]{event} (cod.south);
\draw[inna,->] (cod.south east) -- node[al, pos=0.55]{vibrate} (mot.north);
\draw[tcpa,<->] (wifi.west) -- ++(-1.1,0) node[left, font=\scriptsize, text=tcpc, align=right]{TCP\\and};
\draw[udpa,<->] ([yshift=-3mm]wifi.west) -- ++(-1.1,0) node[left, font=\scriptsize, text=udpc]{UDP};
\end{tikzpicture}
\caption{Work on the two cores of the ESP32. The buffer is the only shared part.}
\end{figure}

\textbf{What is inside the board.} The figure below is a side cut of the board. The motion sensor is
fixed at the centre, under the plywood, so it feels a hit from every side. The four vibration
motors are under the four corners. They have their own power line, so they do not make the LED
voltage drop. The ESP32, the level shifter and the power supply sit under the board too.

\begin{figure}[htbp]\centering
\begin{tikzpicture}[x=0.9cm, y=0.9cm]
\boardside
\nmark{3.3}{-0.15}{1}\nmark{7.0}{0.3}{2}\nmark{1.8}{0.05}{3}\nmark{12.2}{0.05}{3}\nmark{10.1}{-0.15}{4}
\nmark{6.25}{1.7}{5}\nmark{6.25}{4.4}{6}\nmark{12.9}{2.55}{7}
\node[anchor=north, font=\scriptsize] at (7,-0.6) {\begin{tabular}{@{}c@{\ \ }l@{\qquad}c@{\ \ }l@{\qquad}c@{\ \ }l@{}}
\textbf{1} & ESP32 and level shifter & \textbf{4} & 5 V power supply & \textbf{6} & 3D-printed buildings \\
\textbf{2} & Motion sensor at the centre & \textbf{5} & LED strips on the roads & \textbf{7} & Team store with two LEDs \\
\textbf{3} & Vibration motors at the corners & & & & \end{tabular}};
\end{tikzpicture}
\caption{Side cut of the board with the parts that M7 uses.}
\end{figure}

\textbf{Finding a hit.} The sensor measures the motion 1000 times per second. First a filter takes
away gravity and slow tilting, so only the shaking is left. The detector also learns the normal
noise of the room. The start threshold is the mean noise plus six times its spread, so it goes up
on a noisy table. When the signal passes the threshold, a hit starts. The detector takes a 150 ms
window from this moment. The energy E in the window gives the strength with a log scale, like a real
earthquake scale: strength = k × log10(E / E0). Here E0 is the energy of the weakest hit we accept,
and k is set in calibration with a standard hit. The result is cut to 0--10. The first push of the
board shows the direction of the hit, and its size shows how far the hit is from the centre. From
this the detector picks one of 9 zones in a 3 × 3 grid and gives a confidence value. After a hit the
detector waits 1 second, so one hit makes only one event. The event goes to the server within
200 ms.

\begin{figure}[htbp]\centering
\begin{tikzpicture}[ck]
\begin{scope}
\fill[adm!15] (3.0,-1.7) rectangle (4.5,1.9);
\draw[black!40] (-0.2,0) -- (8.2,0);
\draw[black!40,->] (0,-1.7) -- (0,2.0);
\draw[obs, line width=0.9pt, dash pattern=on 4pt off 2pt] (0,0.45) -- (8,0.45);
\draw[obs, line width=0.9pt, dash pattern=on 4pt off 2pt] (0,-0.45) -- (8,-0.45);
\draw[srv, line width=0.8pt, smooth, samples=400, domain=0:8] plot (\x, {0.12*sin(1500*\x) + 0.06*sin(1900*\x) + (\x>3 ? 1.6*exp(-1.6*max(\x-3,0))*sin(2200*(\x-3)) : 0)});
\node[font=\scriptsize, text=adm!80!black] at (3.75,2.15) {150 ms window};
\node[font=\scriptsize, text=obs, anchor=west] at (5.6,0.7) {threshold};
\node[font=\scriptsize, anchor=west] at (0.3,-1.1) {room noise};
\node[font=\scriptsize] at (4.2,-2.0) {time};
\node[al, draw=cli, text=cli!70!black, line width=0.8pt] at (6.9,1.5) {event sent};
\draw[cli, ->, line width=1pt] (4.5,1.5) -- (6.25,1.5);
\end{scope}
\begin{scope}[shift={(9.4,-1.5)}]
\foreach \i in {0,1,2}{\foreach \j in {0,1,2}{%
  \pgfmathtruncatemacro{\z}{(2-\j)*3+\i}
  \filldraw[fill=black!6, draw=black!40] (\i,\j) rectangle ++(1,1);
  \node[font=\tiny, text=black!50] at (\i+0.18,\j+0.82) {\z};}}
\fill[obs!45] (2,1) rectangle ++(1,1);
\node[font=\tiny, text=black!70] at (2.18,1.82) {5};
\fill[brd] (1.5,1.5) circle (0.09);
\draw[obs, line width=1pt, ->] (1.5,1.5) -- (2.5,1.45);
\fill[obs] (2.55,1.45) circle (0.1);
\node[font=\scriptsize] at (1.5,-0.3) {zone 5, confidence 0.8};
\node[font=\scriptsize] at (1.5,3.3) {board from above};
\end{scope}
\end{tikzpicture}
\caption{Left: the sensor signal of one hand hit. Right: the 3 × 3 zones; the purple dot is the sensor.}
\end{figure}

\textbf{Motors.} When the admin presses the earthquake button, M1 sends a vibrate command on TCP.
The command says which motors to use, the strength and the time, at most 5 seconds. M7 drives the
motors with PWM through transistors. It starts slowly over 100 ms, so the current does not jump.
The shaking is then measured by the same detector as a hand hit, and the event is marked with
``motors'' as its source. So the earthquake in the game follows the real shaking of the table.

\textbf{When the link is lost.} If no state comes for 2 seconds, M7 stops the motors at once and
tells M8 to show a ``searching'' animation. Then it tries to connect again. The goal is to be back
within 5 seconds after the router is up. The Wi-Fi name, the password and the calibration values
are stored in the ESP32 flash. We can change them over the USB cable, so we do not need to build
the code again at a new place.

\begin{figure}[htbp]\centering
\begin{tikzpicture}[ck]
\node[st, draw=black!50, fill=black!6] (b) at (0,0) {Start};
\node[st, draw=adm, fill=admL] (w) at (3.4,0) {Join Wi-Fi};
\node[st, draw=adm, fill=admL] (f) at (6.8,0) {Find server};
\node[st, draw=cli, fill=cliL, line width=1.4pt] (c) at (10.2,0) {Connected};
\node[st, draw=obs, fill=obs!10] (s) at (6.8,-1.8) {Searching};
\draw[inna,->] (b) -- (w);
\draw[inna,->] (w) -- node[al]{ok} (f);
\draw[inna,->] (f) -- node[al]{found} (c);
\draw[inna,->] (c) |- node[al, pos=0.3]{no state 2 s} (s);
\draw[inna,->] (s) -- node[al]{retry} (f);
\draw[inna,->] (s) -| node[al, pos=0.75]{no Wi-Fi} (w);
\end{tikzpicture}
\caption{Link states of the board. In ``Searching'' the motors stop and the LEDs show a waiting animation.}
\end{figure}

\msection{Testing and demonstration}

Two PC programs take the place of the neighbours. The server emulator sends test scenarios with
delay and loss that we can change. It also shows the earthquake events and the seismic signal. The
board emulator draws the board on the screen, so M1 and M2 can test before the real board is ready.
The tests use the shared example messages, broken packets, a 2-hour run, added delay, 20\% loss and
a router restart. For the sensor we test hit finding and false alarms. Our targets are: at least
95\% of calibrated hits are found; at most one false earthquake in 30 minutes of normal play; the
correct quarter of the board in at least 95\% of hits and the correct zone in 80\%; and at most
200 ms from the hit to the event. In the demonstration we hit the board by hand and press the admin
button, and we watch the seismic graph and the shock wave on the LEDs.

\end{document}
